\section{Анализ работ}

Работы отобраны по трём вопросам: как карта выражает свою ненадёжность, как планировщик получает гарантию и что происходит с картой при ошибке позы. Все числа ниже взяты из полных текстов (раздел или таблица указаны); площадки подтверждены по DOI или официальной программе, препринты отмечены.

\subsection*{Открытословарные карты и их деградация}

\textbf{ConceptGraphs}~\cite{conceptgraphs} (ICRA~2024) строит из RGB-D-последовательности 3D-граф объектов: маски сегментатора объединяются между кадрами, каждому объекту ставится в соответствие эмбеддинг CLIP, метка --- аргмаксимум косинусного сходства с текстовым запросом. Это де-факто стандарт открытословарного картирования. \emph{Ограничение для навигации}: карта не выдаёт никакой меры уверенности, а аргмаксимум не откалиброван.

\textbf{OSMa-Bench}~\cite{osmabench} (BE2R, IROS~2025) оценивает ConceptGraphs, BBQ и OpenScene на симулированных сценах ReplicaCAD и HM3D при разных условиях освещения; метрики --- mAcc, f-mIoU и ответы на вопросы по графу сцены. Главный результат --- сильная деградация при смене освещения (f-mIoU ConceptGraphs 28{,}00~$\to$~14{,}29). \emph{Ограничение}: бенчмарк не содержит задачи движения, поэтому не говорит, как деградация карты отражается на роботе. \emph{Наблюдение при воспроизведении}: в опубликованном наборе на HuggingFace для каждой сцены есть только одно условие освещения (\texttt{baseline} для 22 сцен ReplicaCAD, \texttt{no\_lights} для 8 сцен HM3D), хотя карточка набора обещает четыре; остальные условия воспроизводятся только генератором HaDaGe на GPU.

\textbf{OSMa-Bench++}~\cite{osmabenchpp} (BE2R, воркшоп ICRA~2026) расширяет протокол на 40 сгенерированных по текстовым описаниям сцен (160 последовательностей, четыре конфигурации освещения) с уклоном в манипуляцию: мелкие объекты, отношения «на», «внутри». На 3~октября 2026~г. данные не опубликованы (ссылка на HuggingFace на странице проекта --- заглушка).

\textbf{On the Overconfidence Problem in Semantic 3D Mapping}~\cite{overconf} (ICRA~2024): наивное байесовское слияние покадровых вероятностей делает карту переуверенной --- средняя ошибка калибровки mECE вокселей 0{,}451 против 0{,}124 у пикселей (табл.~I); усреднение калибрует лучше. Отсюда в нашей карте --- слияние усреднением.

\subsection*{Гарантии для планировщика}

Конформное предсказание пришло в планирование через прогноз движения агентов: Lindemann и др.~\cite{lindemann} калибруют ошибку прогноза траекторий и превращают её в запас безопасности планировщика. Ниже --- работы, где калибруется само восприятие.


\textbf{Sundarsingh и др.}~\cite{sundarsingh} (RA-L~2026) --- ближайшая работа: семантическая карта с конформно откалиброванными множествами меток клеток и планировщик, обходящий клетки, которые \emph{могут} принадлежать запретному классу. Калибровочная единица --- среда (50 калибровочных, 61 тестовая). При $1-\alpha = 95\,\%$ доля корректных карт $93{,}36 \pm 3{,}19\,\%$, успех миссий $98{,}36 \pm 1{,}63\,\%$ (табл.~I). \emph{Ограничения}: поза явно считается известной (разд.~II), классов пять (закрытый словарь), сетка 1~м; при сдвиге условий покрытие падает с $90{,}16$ до $80{,}33\,\%$ при номинале $90\,\%$ (разд.~V-F); в будущей работе авторы называют прямую оценку неопределённости сенсоров.

\textbf{Perceive With Confidence}~\cite{pwc} (IJRR~2025) калибрует \emph{минимальное раздувание} карты занятости, при котором она накрывает истинные препятствия; гарантия не зависит от планировщика. В симуляции 400 калибровочных сред дают при $\varepsilon = 0{,}15$ порог $0{,}75$~м. \emph{Ограничения}: только геометрия, без семантики; состояние робота считается известным (разд.~VI). Наш счёт «расстояние промаха» --- семантический аналог этого счёта.

\textbf{Learnable CP}~\cite{learnablecp} (ICRA~2026) калибрует зазор RRT$^*$ по выученному из 20 геометрических признаков счёту; среди деградаций есть синтетический дрейф локализации ($\sigma = 0{,}5$~м). \emph{Ограничения}: семантики нет, калибровка одна на все среды и виды шума, покрытие --- по точкам пути, а не по задаче.

\textbf{Conformal Constraint Tightening}~\cite{tightening} (препринт 2026, группа Kantaros) калибрует наихудшее отклонение и ужесточает ограничения так, что \emph{любой} план, решающий ужесточённую задачу, безопасен с вероятностью не ниже $1-\alpha$; показано, что доля задач, решаемых RRT, падает с ростом ужесточения. Это прямо задаёт место планировщика: гарантию даёт калибровка, а от планировщика зависит, решит ли он суженную задачу.

\subsection*{Композиция гарантий и ошибка позы}

\textbf{Marginal Calibration Does Not Compose}~\cite{baral} (воркшоп IROS~2026): два модуля, каждый с покрытием $95\,\%$, после композиции дают от $86{,}46$ до $99{,}98\,\%$ в зависимости от корреляции ошибок; при положительной корреляции --- $39{,}7$ столкновения против $8{,}7$ на $10\,000$ прогонов; совместная калибровка возвращает $95{,}18$--$95{,}25\,\%$ (разд.~V, гауссов пример).

\textbf{P-POSEMEM}~\cite{pposemem} (препринт 2026) --- семантическая память, согласованная с перестройкой графа поз SLAM. На реальном роботе замыкания циклов сдвигают ключевые кадры на 2--11~м (медиана), объект с координатой, замороженной при вставке, оказывается в 2{,}75~м от себя, «переключение цели» происходит в $18\,\%$ запросов; при этом номинальный $95\,\%$-эллипсоид позы покрывает лишь $33\,\%$ реальных ошибок. Ошибка позы в семантической навигации --- реальная и плохо откалиброванная.

\subsection*{Сводка}

\begin{table}[H]
\centering\small
\caption{Кто что калибрует и при каком допущении о позе}
\label{tab:lit}
\begin{tabularx}{\linewidth}{L{3.2cm}YYL{2.6cm}}
\toprule
Работа & Что калибруется & Поза & Планировщик в контуре \\
\midrule
Sundarsingh~\cite{sundarsingh} & множества меток клеток (закрытый словарь) & известна & да \\
PwC~\cite{pwc} & раздувание карты занятости & известна & да (любой безопасный) \\
Learnable CP~\cite{learnablecp} & геометрический зазор по признакам & синтетический сдвиг & да (RRT$^*$) \\
Tightening~\cite{tightening} & отклонение исполнения & --- & да (любой) \\
P-POSEMEM~\cite{pposemem} & --- (байесовская память) & распределение, некалибровано & цели, без гарантии \\
\textbf{Эта работа} & \textbf{расстояние промаха: метка + поза} & \textbf{оценка с дрейфом} & \textbf{да} \\
\bottomrule
\end{tabularx}
\end{table}

Пустая клетка таблицы~\ref{tab:lit} --- одна distribution-free гарантия на ошибку метки открытого словаря \emph{и} ошибку позы, проверенная с планировщиком в контуре. По цитированиям ближайших работ и поиску по arXiv и OpenAlex (2~октября 2026~г.) такой работы не найдено; поиск Semantic Scholar по ключевым словам в эти дни не работал, поэтому вывод о новизне предварительный.
